% =============================================================
%  abstract.tex -- text of the dissertation Abstract (<= 300 words)
%
%  Draft 2026-10-29. The environment is NOT in this file: dissertation.tex
%  already has \begin{abstract} ... \end{abstract}, so replace the
%  \todo{Write abstract.} line inside it with
%      % =============================================================
%  abstract.tex -- text of the dissertation Abstract (<= 300 words)
%
%  Draft 2026-10-29. The environment is NOT in this file: dissertation.tex
%  already has \begin{abstract} ... \end{abstract}, so replace the
%  \todo{Write abstract.} line inside it with
%      % =============================================================
%  abstract.tex -- text of the dissertation Abstract (<= 300 words)
%
%  Draft 2026-10-29. The environment is NOT in this file: dissertation.tex
%  already has \begin{abstract} ... \end{abstract}, so replace the
%  \todo{Write abstract.} line inside it with
%      % =============================================================
%  abstract.tex -- text of the dissertation Abstract (<= 300 words)
%
%  Draft 2026-10-29. The environment is NOT in this file: dissertation.tex
%  already has \begin{abstract} ... \end{abstract}, so replace the
%  \todo{Write abstract.} line inside it with
%      \input{Introduction/abstract}
%
%  Numbers are written out (not macros) because the abstract is typeset
%  before the chapter value files are \input. Values come from
%  Chapter2/Writing/chapter2_derived_values.tex (Ea = 157.1 +/- 5.3 kJ/mol)
%  and Chapter3's conclusions (cumulative yields <= 5% up to 600 C).
%
%  NOTE: Do NOT include \documentclass or \begin{document} here.
% =============================================================

Rising demand for critical metals, driven by computing infrastructure and the
transition to renewable energy, foreground the need to understand how ore deposits form,
change, and survive. This thesis develops geochemical tools for three families of
processes that concentrate and rework rare metals: magmatic, metasomatic, and
meteoric. The first three chapters extend U-Th/He geochronology to minerals that
decompose under the vacuum of conventional helium step-heating experiments.
Chapter~\ref{ch:apparatus} presents a method that seals samples in helium-permeable
silica ampoules holding water- or sulfur-saturated conditions, with a framework for
disentangling the proceeses of helium permeation through the glass from release by the sample.
Chapter~\ref{ch:goethite} applies this method to goethite, a direct chronometer of weathering
and laterite formation. Under water-saturated conditions helium release is diffusive,
with an activation energy of $157 \pm 5$~kJ/mol and near-quantitative retention at
surface temperatures over billion-year timescales for the sample studied, whereas
earlier vacuum experiments recorded conversion to hematite.
Chapter~\ref{ch:sulfide} shows that pyrite and arsenopyrite are highly helium retentive, with pyrite in particular being among the most
helium-retentive minerals studied, comparable to or exceeding monazite. 
U-Th/He ages from the Gympie gold field and the Kidd Creek
massive sulfide deposit depend on uranium- and thorium-bearing inclusions and must be
interpreted locality by locality.
Chapter~\ref{ch:ngibbs} presents \emph{nGibbs}, a software package of compositionally general deep-learning emulators of the
\textsc{MELTS} and \textsc{HeFESTo} thermodynamic models that conserve mass and run
two to four orders of magnitude faster, making composition a practical free variable
in models of magmatic and mantle systems.
Chapter~\ref{ch:draaikraal} turns to the Bushveld Complex, the world's dominant source
of platinum-group metals. Field mapping at Draaikraal, Mpumalanga shows that local deformation at the 
intrusion's floor is consistent with a synmagmatic origin, and that the area's recessive,
laterite-capped topography is the result of contamination by fluids from devolatilizing floor
sediments. The Bushveld, with its sulfide-hosted platinum-group metals, goethite
laterites, and unresolved petrological puzzles, is a natural place to apply these tools
together.

%
%  Numbers are written out (not macros) because the abstract is typeset
%  before the chapter value files are \input. Values come from
%  Chapter2/Writing/chapter2_derived_values.tex (Ea = 157.1 +/- 5.3 kJ/mol)
%  and Chapter3's conclusions (cumulative yields <= 5% up to 600 C).
%
%  NOTE: Do NOT include \documentclass or \begin{document} here.
% =============================================================

Rising demand for critical metals, driven by computing infrastructure and the
transition to renewable energy, foreground the need to understand how ore deposits form,
change, and survive. This thesis develops geochemical tools for three families of
processes that concentrate and rework rare metals: magmatic, metasomatic, and
meteoric. The first three chapters extend U-Th/He geochronology to minerals that
decompose under the vacuum of conventional helium step-heating experiments.
Chapter~\ref{ch:apparatus} presents a method that seals samples in helium-permeable
silica ampoules holding water- or sulfur-saturated conditions, with a framework for
disentangling the proceeses of helium permeation through the glass from release by the sample.
Chapter~\ref{ch:goethite} applies this method to goethite, a direct chronometer of weathering
and laterite formation. Under water-saturated conditions helium release is diffusive,
with an activation energy of $157 \pm 5$~kJ/mol and near-quantitative retention at
surface temperatures over billion-year timescales for the sample studied, whereas
earlier vacuum experiments recorded conversion to hematite.
Chapter~\ref{ch:sulfide} shows that pyrite and arsenopyrite are highly helium retentive, with pyrite in particular being among the most
helium-retentive minerals studied, comparable to or exceeding monazite. 
U-Th/He ages from the Gympie gold field and the Kidd Creek
massive sulfide deposit depend on uranium- and thorium-bearing inclusions and must be
interpreted locality by locality.
Chapter~\ref{ch:ngibbs} presents \emph{nGibbs}, a software package of compositionally general deep-learning emulators of the
\textsc{MELTS} and \textsc{HeFESTo} thermodynamic models that conserve mass and run
two to four orders of magnitude faster, making composition a practical free variable
in models of magmatic and mantle systems.
Chapter~\ref{ch:draaikraal} turns to the Bushveld Complex, the world's dominant source
of platinum-group metals. Field mapping at Draaikraal, Mpumalanga shows that local deformation at the 
intrusion's floor is consistent with a synmagmatic origin, and that the area's recessive,
laterite-capped topography is the result of contamination by fluids from devolatilizing floor
sediments. The Bushveld, with its sulfide-hosted platinum-group metals, goethite
laterites, and unresolved petrological puzzles, is a natural place to apply these tools
together.

%
%  Numbers are written out (not macros) because the abstract is typeset
%  before the chapter value files are \input. Values come from
%  Chapter2/Writing/chapter2_derived_values.tex (Ea = 157.1 +/- 5.3 kJ/mol)
%  and Chapter3's conclusions (cumulative yields <= 5% up to 600 C).
%
%  NOTE: Do NOT include \documentclass or \begin{document} here.
% =============================================================

Rising demand for critical metals, driven by computing infrastructure and the
transition to renewable energy, foreground the need to understand how ore deposits form,
change, and survive. This thesis develops geochemical tools for three families of
processes that concentrate and rework rare metals: magmatic, metasomatic, and
meteoric. The first three chapters extend U-Th/He geochronology to minerals that
decompose under the vacuum of conventional helium step-heating experiments.
Chapter~\ref{ch:apparatus} presents a method that seals samples in helium-permeable
silica ampoules holding water- or sulfur-saturated conditions, with a framework for
disentangling the proceeses of helium permeation through the glass from release by the sample.
Chapter~\ref{ch:goethite} applies this method to goethite, a direct chronometer of weathering
and laterite formation. Under water-saturated conditions helium release is diffusive,
with an activation energy of $157 \pm 5$~kJ/mol and near-quantitative retention at
surface temperatures over billion-year timescales for the sample studied, whereas
earlier vacuum experiments recorded conversion to hematite.
Chapter~\ref{ch:sulfide} shows that pyrite and arsenopyrite are highly helium retentive, with pyrite in particular being among the most
helium-retentive minerals studied, comparable to or exceeding monazite. 
U-Th/He ages from the Gympie gold field and the Kidd Creek
massive sulfide deposit depend on uranium- and thorium-bearing inclusions and must be
interpreted locality by locality.
Chapter~\ref{ch:ngibbs} presents \emph{nGibbs}, a software package of compositionally general deep-learning emulators of the
\textsc{MELTS} and \textsc{HeFESTo} thermodynamic models that conserve mass and run
two to four orders of magnitude faster, making composition a practical free variable
in models of magmatic and mantle systems.
Chapter~\ref{ch:draaikraal} turns to the Bushveld Complex, the world's dominant source
of platinum-group metals. Field mapping at Draaikraal, Mpumalanga shows that local deformation at the 
intrusion's floor is consistent with a synmagmatic origin, and that the area's recessive,
laterite-capped topography is the result of contamination by fluids from devolatilizing floor
sediments. The Bushveld, with its sulfide-hosted platinum-group metals, goethite
laterites, and unresolved petrological puzzles, is a natural place to apply these tools
together.

%
%  Numbers are written out (not macros) because the abstract is typeset
%  before the chapter value files are \input. Values come from
%  Chapter2/Writing/chapter2_derived_values.tex (Ea = 157.1 +/- 5.3 kJ/mol)
%  and Chapter3's conclusions (cumulative yields <= 5% up to 600 C).
%
%  NOTE: Do NOT include \documentclass or \begin{document} here.
% =============================================================

Rising demand for critical metals, driven by computing infrastructure and the
transition to renewable energy, foreground the need to understand how ore deposits form,
change, and survive. This thesis develops geochemical tools for three families of
processes that concentrate and rework rare metals: magmatic, metasomatic, and
meteoric. The first three chapters extend U-Th/He geochronology to minerals that
decompose under the vacuum of conventional helium step-heating experiments.
Chapter~\ref{ch:apparatus} presents a method that seals samples in helium-permeable
silica ampoules holding water- or sulfur-saturated conditions, with a framework for
disentangling the proceeses of helium permeation through the glass from release by the sample.
Chapter~\ref{ch:goethite} applies this method to goethite, a direct chronometer of weathering
and laterite formation. Under water-saturated conditions helium release is diffusive,
with an activation energy of $157 \pm 5$~kJ/mol and near-quantitative retention at
surface temperatures over billion-year timescales for the sample studied, whereas
earlier vacuum experiments recorded conversion to hematite.
Chapter~\ref{ch:sulfide} shows that pyrite and arsenopyrite are highly helium retentive, with pyrite in particular being among the most
helium-retentive minerals studied, comparable to or exceeding monazite. 
U-Th/He ages from the Gympie gold field and the Kidd Creek
massive sulfide deposit depend on uranium- and thorium-bearing inclusions and must be
interpreted locality by locality.
Chapter~\ref{ch:ngibbs} presents \emph{nGibbs}, a software package of compositionally general deep-learning emulators of the
\textsc{MELTS} and \textsc{HeFESTo} thermodynamic models that conserve mass and run
two to four orders of magnitude faster, making composition a practical free variable
in models of magmatic and mantle systems.
Chapter~\ref{ch:draaikraal} turns to the Bushveld Complex, the world's dominant source
of platinum-group metals. Field mapping at Draaikraal, Mpumalanga shows that local deformation at the 
intrusion's floor is consistent with a synmagmatic origin, and that the area's recessive,
laterite-capped topography is the result of contamination by fluids from devolatilizing floor
sediments. The Bushveld, with its sulfide-hosted platinum-group metals, goethite
laterites, and unresolved petrological puzzles, is a natural place to apply these tools
together.
